% =============================================================================
% 1.4 论文组织结构
% 草稿版本：v0.1 (2026-05-05)
% 字数：约 800 字
% 风格规则：v0.2 风格（无 \textbf / \textit / 引例括号 / 双引号 / 破折号）
% 引用：暂无新增引用
% 依赖宏包：booktabs / tabularx（主稿已加载）
% =============================================================================

\subsection{论文组织结构}
\label{sec:thesis-structure}

本文按问题分析、方法实现、实验验证三段式逻辑组织内容，共划分为
五章。

第 1 章为绪论。本章给出研究背景与意义，综述气象 AI、大语言模型
智能体技术、检索增强生成与工具学习三个方向的国内外研究现状，并
据此明确本文的研究问题与五项创新贡献，最后说明全文组织结构。

第 2 章为气象预报任务技术问题分析与解决方案。本章把基于 AI 智能体
的气象数据检索与分析任务按数据流向建模为指令解析、异构数据、数值
与逻辑幻觉三个相互衔接的研究问题，并给出基于 ReAct 范式的单智能体
思维链模型作为统一的方法论框架。本章同时给出基于 RAG 的领域知识
注入与基于 Prompt 工程的工具调用约束作为方法论框架在领域上的落地
策略。

第 3 章为基于工具学习的异构气象数据智能检索技术。本章针对指令
解析问题展开，从工具学习、意图识别、参数补全三个维度构建从自然
语言查询到结构化检索请求的完整链路。本章把和风天气与 Open-Meteo
两类异构 API 统一封装为 Function Calling 工具，把自然语言查询
通过 \texttt{recognizer} 与 \texttt{completer} 双阶段管线解析为
\texttt{WeatherIntent} 结构化对象，分层架构本身吸收了单点大语言
模型抖动的容错负担。

第 4 章为融合语义桥接机制的气象数据统计分析方法。本章针对异构数据
与数值幻觉两个研究问题展开，给出双通路 RAG 架构、代码生成加受限
沙箱执行机制、自适应气象决策报告生成三类核心方法。其中通路 A 承担
混合检索召回任务，通路 B 承担确定性分级与硬链接式语义桥接任务,
两条通路共同覆盖概念性查询与数值归属两类需求；代码生成与沙箱执行
解决数值幻觉问题；多场景报告模板把数值数据、语义解读、统计结果
整合为面向终端用户的结构化决策报告。

第 5 章为实验验证与总结。本章先给出实验环境与 5 套评测集的统一
说明，再按意图识别、通路 A 检索、通路 B 桥接、代码执行 4 个模块
层评测加 1 个 30 条端到端集成层评测的顺序展开 5 类系统化实证,
最后从工作总结与研究前景两个维度收束全文。本章同时给出本文方法
在多轮长程依赖、数值阈值精确归属、沙箱安全三方面的局限性回顾,
并提出显式对话状态跟踪、知识库扩展、安全沙箱升级、多智能体协作
四个方向的未来工作展望。

各章之间的逻辑关系可以归纳为问题驱动加方法落地加实证收束的三段
结构。第 2 章是问题与方法的定义层，第 3 至第 4 章是方法的实现层,
第 5 章是实证的验证层。指令解析问题在第 3 章落地，异构数据与数值
幻觉两个问题在第 4 章落地，所有方法的有效性由第 5 章的 5 类系统化
评测共同验证。本文的核心论证链路据此构成完整闭环。

